% results.tex

We evaluated SA on four protein families spanning a range of sizes, functions, and PCA-space
geometries (Table~\ref{tab:cross-family}). Three Pfam seed families (Kunitz domains (PF00014,
$K=99$), SH3 domains (PF00018, $K=55$), and WW domains (PF00397, $K=420$)) were used to
benchmark both hard curation and multiplicity-weighted conditioning, with designated subsets
defined by known binding-site residues. A fourth family, the $\omega$-conotoxin O-superfamily
($K=74$, curated from SwissProt), was used to test hard curation on a therapeutically relevant
peptide system targeting the Cav2.2 calcium channel. For each family, we split sequences into
designated and background subsets based on an experimentally validated phenotype marker,
built conditioned memory matrices, and generated sequences via Langevin dynamics at the
entropy-inflection inverse temperature $\beta^{*}$. The Fisher separation index $S$, a
measure of how well PCA separates designated from background sequences, varied from 0.11
(WW, heavily interleaved) to 0.34 (SH3, well-separated), providing a natural axis along which
to test the limits of multiplicity-weighted conditioning.

\paragraph{Hard curation and the small-set amplifier.}
We restricted the memory matrix to a user-designated functional subset of Kunitz domains
(PF00014; $K=99$ sequences, $L=53$ aligned positions) and tested whether subset-specific
phenotype transfers to generated sequences (Fig.~\ref{fig:scaling}). Kunitz domains are serine protease inhibitors whose binding
specificity depends on the P1 residue at the primary contact site with the target protease: Lys
or Arg at P1 confers strong trypsin inhibition, while other residues (Gly, Gln, Asn) do
not~\cite{laskowskiProteinInhibitors1980}. We designated the 32 sequences with K/R at P1 as the
functional subset ($K_{\mathrm{des}} = 32$) and the remaining 67 as background
($K_{\mathrm{bg}} = 67$), built separate memory matrices from each group, found $\beta^{*}$
for each via entropy inflection, and generated 930 sequences per condition (30 chains $\times$
5{,}000 steps, thinned after burn-in). We observed that 100\% of
designated-subset-conditioned sequences had K/R at P1; 0\% of background-conditioned sequences
did; and full-family generation produced 41\% K/R, close to the natural proportion
($32/99 = 32\%$), which demonstrates that the designated phenotype was
inherited perfectly while the fold-level composition was approximately preserved.
Sequence-level inspection confirmed that the top strong-binder-conditioned sequences
preserved all highly conserved positions ($11/11$), all six cysteines, and carried K/R
at P1, while introducing 19-26 substitutions concentrated at variable sites;
per-position Shannon entropy tracked the stored MSA closely ($r = 0.988$ for full-family,
$r = 0.932$ for strong-conditioned), with the largest deviation at P1 where conditioning
collapsed diversity to K/R (Fig.~\ref{fig:sequence-analysis}).
To determine the minimum input size required for reliable transfer,
we subsampled $K_{\mathrm{des}} \in \{3, 5, 8, 10, 15, 20, 25, 32\}$ from the Kunitz K/R
subset (3~replicates each), built curated memory matrices, and generated sequences. We found
that P1 phenotype fidelity was 100\% at \emph{every} sample size, including
$K_{\mathrm{des}} = 3$, but that small input sets reduced diversity: pairwise sequence diversity
increased from 0.36 ($K_{\mathrm{des}} = 3$) to 0.50 ($K_{\mathrm{des}} = 32$), and KL
divergence decreased from 0.047 to 0.009 over the same range, which is consistent with fewer
stored patterns constraining the PCA subspace and reducing the volume accessible to the
Langevin chain (Fig.~\ref{fig:scaling}). These results validate the ``small-set amplifier''
framing: a practitioner with as few as 5-10 experimentally characterized sequences (for
instance, hits from a phage display screen) can generate hundreds of novel candidates that
maintain the input phenotype while introducing substantial sequence diversity (pairwise diversity
$\geq 0.45$ for $K_{\mathrm{des}} \geq 10$).

\paragraph{Multiplicity-weighted conditioning and the calibration gap.}
We swept $\rho$ from 1 (standard SA) to 1{,}000 on the full Kunitz family ($K = 99$) to test
multiplicity-weighted generation (Fig.~\ref{fig:phase-transition}, Table~\ref{tab:per-family-rho}). For each $\rho$, we
computed the multiplicity vector $\vr$, found $\beta^{*}(\vr)$ via entropy inflection on the
weighted attention distribution, and generated 930 sequences (30 chains). We observed that the
attention weight on designated patterns tracked $f_{\mathrm{eff}}$ with negligible deviation
($< 0.3\%$) at all $\rho$ values, which confirms Proposition~\ref{prop:score}: the multiplicity
weights tilted the energy landscape exactly as intended. The observed P1~K/R fraction increased
monotonically from 0.41 ($\rho=1$) to 0.63 ($\rho=500$), substantially exceeding the natural
proportion but falling short of the $f_{\mathrm{eff}} = 0.996$ predicted by the energy-level
conditioning; sequence diversity decreased modestly from 0.56 to 0.51 over this range and KL
divergence increased from 0.007 to 0.017. The phase transition
$\beta^{*}$ shifted from 4.3 at $\rho=1$ ($K_{\mathrm{eff}} = 99$) to 10.1 at
$\rho=1{,}000$ ($K_{\mathrm{eff}} = 32.1$), and the entropy curves at different $\rho$ values
formed a family of sigmoidal profiles with systematically rightward-shifted inflection points,
which confirms that $\beta^{*}(\vr)$ depends on $K_{\mathrm{eff}}(\vr)$ rather than the raw
pattern count $K$ (Fig.~\ref{fig:phase-transition}). To understand why $f_{\mathrm{obs}}$
lagged $f_{\mathrm{eff}}$, we measured all three layers of the signal chain
(Eq.~\ref{eq:gap-decomp}) independently across a fine $\rho$ sweep ($\rho \in [1, 1000]$,
17~values, 5~replicates each), recording at each step: (i)~the total softmax attention weight
on designated patterns ($\bar{a}_{\mathrm{des}}$), (ii)~the continuous K/R probability at P1 in
the PCA-reconstructed one-hot vector ($f_{\mathrm{soft}}$), and (iii)~the hard-decoded residue
identity ($f_{\mathrm{obs}}$). We found that the attention gap was
$\Delta_{\mathrm{attn}} \approx 0$ across all $\rho$ values ($< 0.3\%$), confirming that the
energy-level conditioning was exact; the PCA gap was the dominant component, with
$f_{\mathrm{soft}}$ effectively flat at $\sim\!0.27$ regardless of $\rho$, varying by less than
0.04 across a thousandfold range of multiplicity ratios; and the argmax gap was negative
($\approx -0.25$), indicating that discrete decoding partially recovered signal lost in the PCA
reconstruction. These observations suggest that the PCA
encoding mapped designated and background sequences to overlapping regions at the P1 position:
even when 99.8\% of attention weight was on the designated patterns, the continuous one-hot
reconstruction could not resolve K~vs.~G at P1 because this variation was not well-captured by
the top $d=80$ principal components. The Fisher separation index for the Kunitz
designated/background split was $S = 0.20$, confirming moderate but incomplete geometric
separation.

\paragraph{Cross-family validation and the separation index.}
We repeated the multiplicity analysis on two additional families (SH3 domains (PF00018;
$K=55$, split by conserved Trp at the binding groove) and WW domains (PF00397; $K=420$, split
by specificity-loop residue)), yielding three families with distinct PCA-space geometries
(Table~\ref{tab:cross-family}, Table~\ref{tab:per-family-rho}). We found that the calibration gap decreased monotonically with
increasing separation index: WW domains ($S = 0.11$, functional subsets heavily interleaved in
PCA space) exhibited the largest gap ($\Delta = 0.63$); Kunitz domains ($S = 0.20$) an
intermediate gap ($\Delta = 0.37$); and SH3 domains ($S = 0.34$, well-separated functional
subsets) a near-zero gap ($\Delta = 0.01$, $f_{\mathrm{obs}} = 0.989$ at
$f_{\mathrm{eff}} = 0.999$). A linear fit yielded $\Delta \approx 0.92 - 2.7\,S$ with
$R^{2} \approx 1.0$ (Fig.~\ref{fig:separation-vs-gap}). Extending the multiplicity sweep to
$\omega$-conotoxin ($S = 0.78$, pharmacologically defined binder split) yielded
$\Delta = 0.15$, following the monotonic trend but lying above the Pfam regression line--consistent
with its imperfect single-residue marker (Tyr13 frequency $= 0.46$ in the full family) and
non-Pfam provenance. We observed that the attention weights tracked $f_{\mathrm{eff}}$
perfectly in all four families, which confirms that the multiplicity-weighted energy is exact
regardless of the downstream calibration gap, while hard curation achieved 100\% phenotype
transfer in all families regardless of separation index--as expected, since curated
memory matrices are built from designated-subset-only PCA, which by construction captures
subset-specific variation in its leading components. To test
whether increasing $\beta$ beyond $\beta^{*}$ could push more phenotype signal through the PCA
bottleneck, we swept $\beta/\beta^{*}$ from 0.5 to 3.0 at three fixed $\rho$ values
($\rho = 10, 50, 200$) and observed a consistent tradeoff: higher $\beta$ increased P1~K/R
fraction at the cost of reduced diversity--at $\rho = 200$ and $\beta = 3\beta^{*}$, P1~K/R
reached 0.71 (versus 0.61 at $\beta^{*}$) while diversity dropped from 0.53 to 0.46--which
suggests that the practitioner can tune $(\rho,\, \beta/\beta^{*})$ jointly to navigate the
fidelity--diversity Pareto frontier.

\paragraph{Application to $\omega$-conotoxin Cav2.2 binders.}
We applied SA to the $\omega$-conotoxin O-superfamily (disulfide-rich peptides from predatory
cone snails that block voltage-gated N-type calcium channels (Cav2.2) in sensory
neurons~\cite{omegaConotoxinSAR2006}) to test whether curated memory conditioning generalizes
beyond classical protein domains to therapeutically relevant peptide families
(Table~\ref{tab:conotoxin-pharmacophore}, Fig.~\ref{fig:conotoxin-loop}). Cav2.2 is a primary
gatekeeper of pain signaling in dorsal root ganglion (DRG) neurons: channel block attenuates
synaptic transmission between peripheral nociceptors and the spinal cord. This family has direct
therapeutic relevance--$\omega$-conotoxin MVIIA (ziconotide) is an FDA-approved non-opioid
analgesic for intractable pain~\cite{ziconotideReview2004}--and is motivated by prior
computational modeling of ATP-gated calcium dynamics in DRG neurons that identified Cav2.2 as a
central node in the pain signaling network~\cite{songModelingAnalysis2009}. The primary Cav2.2
pharmacophore is Tyr13 (MVIIA numbering): the hydroxyl group contacts the channel pore directly
and its substitution abolishes binding, while basic residues (Lys, Arg) distributed across the
peptide contribute electrostatic complementarity with the negatively charged channel entrance.
We compiled 74 SwissProt $\omega$-conotoxin sequences, computed a multiple sequence alignment
with MAFFT~\cite{mafft2013}, retained 26 positions after gap filtering, designated 23 sequences
as strong Cav2.2 binders (confirmed or predicted N-type channel blockers, including MVIIA,
GVIA, CVID, CVIE, CVIF, and close homologs) and 2 sequences as weak or non-N-type-selective
(MVIIC and MVIID, which preferentially block P/Q-type channels with
IC$_{50} \geq 200$\,nM against Cav2.2), and ran two independent SA experiments (50~chains
$\times$ 5\,000 steps, $\alpha = 0.01$), generating 1\,550 sequences per condition: one seeded
from the full family ($K = 74$, $\beta^{*} = 3.23$) and one from the strong binders only
($K = 23$, $\beta^{*} = 2.72$). We found that pharmacophore inheritance was near-complete under
strong-binder seeding and substantially attenuated under full-family seeding: in the full-family
input only 33.8\% of the 74 sequences carried Tyr at the pharmacophore position (col.~13), and
full-family-seeded generation raised this modestly to 46.9\%, whereas 82.6\% of the 23 strong
binders carried Tyr13 and strong-binder-seeded SA amplified this to 98.3\%--essentially
locking the primary Cav2.2 pharmacophore across the generated library
(Table~\ref{tab:conotoxin-pharmacophore}). The basic-residue (Lys + Arg) fraction was also
faithfully inherited: strong-binder generation produced 20.1\% basic residues versus 19.7\% in
the input, while full-family generation produced only 12.0\%, matching the lower K/R content of
the diverse full-family input. Per-position amino acid frequency heatmaps confirmed the
conditioning effect visually (Fig.~\ref{fig:conotoxin-loop}): the strong-binder input and strong-binder-generated heatmaps were
nearly indistinguishable; deep-red Tyr at position~13, conserved Cys framework at
positions~14-16, variable-loop diversity at positions~9-12; and the residual heatmaps
(generated $-$ input) showed near-zero deviations across all positions, which suggests that SA
faithfully reproduced the input frequency statistics rather than collapsing to a consensus
sequence. Sequence-level inspection of the five
highest-confidence strong-seeded sequences (ranked by ESMFold pLDDT) confirmed that all five
preserved every highly conserved position ($6/6$), all six cysteines, and carried Tyr at the
pharmacophore position, while introducing 5-11 substitutions at variable sites; the
per-position entropy correlation between stored and strong-seeded sequences was $r = 0.745$,
lower than the near-perfect full-seeded correlation ($r = 0.997$), reflecting the collapse of
diversity at co-varying loop positions that contribute to Cav2.2 binding specificity
(Fig.~\ref{fig:conotoxin-sequence-analysis}). We further validated the generated sequences
against published structure--activity relationship (SAR) data by measuring agreement with
experimentally identified binding determinants at 12 positions
(Table~\ref{tab:sar-agreement}): SA strong-seeded sequences preserved the primary
pharmacophore (Tyr13: 98.3\%), all six disulfide-forming cysteines
($\geq 97.4\%$), and enriched basic residues at loop~2 positions critical for channel
interaction (Lys2: 97.5\% K/R; Arg10: 82.3\% K/R) relative to both the 23-sequence input
and full-family generation, which demonstrates that curated memory conditioning not only
transferred global pharmacophore statistics but also preserved the specific residue-level
binding determinants identified by decades of mutagenesis
work~\cite{kimTyr13Essential1995,satoBasicResidues1993,lewisConotoxinSAR2012}.
Structural validation confirmed that the generated conotoxin sequences adopted the
expected inhibitor cystine knot fold, with the characteristic disulfide-constrained loop
architecture visually preserved across all three top SA strong-seeded variants
(Fig.~\ref{fig:fold-superposition}). Together,
these results confirm that curated memory conditioning transferred Cav2.2 pharmacophore
statistics from a small set of 23 characterized sequences to a diverse library of 1\,550
novel candidates, with no training and no target structure required.

\paragraph{AlphaFold2 multimer binding validation.}
We assessed whether preserved pharmacophores translate to predicted binding competence by performing AlphaFold2 multimer complex prediction via ColabFold for representative SA-generated sequences paired with the Cav2.2 pore domain target (Table~\ref{tab:docking-validation}). We submitted 9 SA strong-binder-conditioned sequences, 4 SA full-family-conditioned sequences, and 5 natural $\omega$-conotoxin controls (CVIE, CVIF, CnVIIE, CnVIIF, CnVIIG) and measured complex confidence using the predicted Template Modeling (pTM) score. We found that SA strong-binder-conditioned sequences achieved complex confidence scores statistically equivalent to natural controls (0.377±0.011 vs 0.362±0.029; Welch t-test: $t=1.070$, $p>0.05$), as did SA full-family-conditioned sequences (0.355±0.019 vs controls; $t=-0.429$, $p>0.05$). Mann-Whitney U tests confirmed the absence of significant differences (all $p>0.05$), which demonstrates that pharmacophore inheritance successfully translates to predicted three-dimensional binding capability. These results provide computational validation that SA conditioning produces $\omega$-conotoxin candidates with predicted Cav2.2 binding competence equivalent to natural channel blockers.

\paragraph{Structural and linguistic validation.}
We predicted folds for 50 sequences from each source (stored, SA full-family, SA
strong-conditioned, SA weak-conditioned, and HMM baseline) using ESMFold and computed TM-scores
against the experimental Kunitz domain structure (1BPI, chain~A) to verify that SA-generated
sequences encode plausible three-dimensional structures
(Table~\ref{tab:structure-validation}). We found that
SA-generated sequences were structurally indistinguishable from natural sequences: mean pLDDT
was $90.4 \pm 1.2$ (SA full) and $90.7 \pm 1.1$ (SA strong) versus $89.3 \pm 3.2$ for stored
sequences, with 100\% of SA sequences exceeding the pLDDT $>70$ confidence threshold; mean
TM-scores were $0.84 \pm 0.01$ (SA full) and $0.83 \pm 0.01$ (SA strong) versus
$0.83 \pm 0.03$ for stored sequences, all well above the $\mathrm{TM} > 0.5$ same-fold
threshold (Fig.~\ref{fig:fold-superposition}); whereas HMM-emitted sequences achieved only $60.4 \pm 13.1$ mean pLDDT (20\% above
threshold) and $0.60 \pm 0.10$ mean TM-score, which suggests that profile HMM emission does not
reliably produce foldable sequences (Table~\ref{tab:structure-validation}). AlphaFold2-ptm predictions via ColabFold~\cite{jumperHighlyAccurate2021,mirdita2022colabfold} independently confirmed these results (Table~\ref{tab:af2-validation}): Kunitz TM-scores were $0.83$--$0.85$ (pLDDT $93$--$95$) and $\omega$-conotoxin TM-scores were $0.35$--$0.47$ (pLDDT $68$--$79$) across stored and SA-generated sequences, consistent with the ESMFold values and the shorter peptide length of conotoxins. We further assessed
sequence plausibility using ESM2-650M pseudo-perplexity computed via masked marginal scoring
over all positions and found that SA-generated sequences achieved pseudo-perplexity comparable
to natural sequences (SA full: $4.78 \pm 0.64$; SA strong: $4.72 \pm 0.60$; stored:
$4.97 \pm 0.98$), while HMM-emitted sequences were approximately threefold worse
($16.1 \pm 4.2$)--the fact that conditioned SA sequences were slightly \emph{more} plausible
than stored sequences is consistent with SA's tendency to generate near the center of the family
distribution, avoiding outlier sequences present in natural alignments
(Table~\ref{tab:structure-validation}). We also compared SA against HMMER3~\cite{hmmer3} profile HMM emission
(150 sequences, \texttt{hmmemit --seed 42}) and bootstrap resampling from stored patterns and
found that SA full-family generation preserved the natural P1~K/R proportion
($0.41 \pm 0.01$) and achieved the lowest KL divergence ($0.0069 \pm 0.0001$), while HMM
emission produced sequences with substantially lower P1~K/R fraction ($0.15 \pm 0.03$) and
higher compositional divergence ($0.030 \pm 0.011$); bootstrap resampling trivially preserved
composition (KL $\approx 0.001$) but produced no novel sequences; and SA strong-binder
conditioning achieved perfect P1~K/R fraction ($1.0 \pm 0.0$)--a capability fundamentally
unavailable to unconditional baselines (Table~\ref{tab:structure-validation}). These results
confirm that SA generated structurally valid, language-model-plausible sequences that
outperformed classical baselines on both fidelity and functional conditioning.
